\chapter{Symbolic Integration}
\label{ch:integration}

% Written in the v0.507 -> v0.6664 book pass (BOOK_UPDATE_REVIEW.md I1-I6, R18).

Differentiation is a finite recursion. Apply the sum, product and chain rules, and
any formula built from the elementary functions has a derivative of the same kind.
Integration runs the other way and is a different kind of problem. There is no
chain rule in reverse. The antiderivative of a short formula can be long, can
involve logarithms of expressions that appear nowhere in the integrand, or can fail
to exist among the elementary functions at all. For two centuries the subject was
a catalogue of tricks. The Risch algorithm and its successors replaced the
catalogue with a theorem. For a large and precisely defined class of integrands,
the question ``does this function have an elementary antiderivative, and if so,
what is it?'' can be \emph{decided} by a finite computation.

\B{Integrate} is built on that theory. This chapter follows it from the easiest
case to the hardest. Section~\ref{sec:int-shape} states the problem precisely and
gives Liouville's theorem, which makes it decidable.
Section~\ref{sec:int-rational} integrates rational functions, the case every
other case reduces to. Section~\ref{sec:int-tower} builds the transcendental
\emph{tower} of logarithms, exponentials, tangents and Liouvillian special
functions, and decides integrability in it with the Risch differential equation.
Section~\ref{sec:int-trager} turns to algebraic functions, where the answer lives
on an algebraic curve and Trager's algorithm decides it.
Section~\ref{sec:int-bronstein} combines the two: a radical over a tower of
logarithms and exponentials, in Bronstein's algorithm.
Section~\ref{sec:int-pmt} describes the \texttt{"ParallelMixedTower"} method, a
parallel Risch--Norman search that reaches integrands outside the decided
classes. Section~\ref{sec:int-cascade} puts the methods together in the order
\B{Integrate} tries them, explains what its messages promise, compares the result
with FriCAS, and lists what is still open.

One promise runs through the chapter. When \B{Integrate} returns an
antiderivative, it has differentiated it and checked that the derivative equals
the integrand. When it prints \texttt{Integrate::nonelem}, it holds a proof that no
elementary antiderivative exists. When it cannot decide, it says so with
\texttt{Integrate::incomplete}, or returns the integral unevaluated. It does not
guess.

%======================================================================
\section{The shape of the problem}
\label{sec:int-shape}

Start with what \B{Integrate} returns on a few familiar integrands.

\begin{codepairs}
\pair{integration/integrate-special/1}{Integration by parts, done for you:
\(\int x e^x = (x-1)e^x\).}
\pair{integration/integrate-special/2}{\(e^{-x^2}\) has no elementary integral, so
the answer is the error function \B{Erf}.}
\pair{integration/integrate-special/3}{\(e^x/x\) gives the exponential integral
\B{ExpIntegralEi}.}
\pair{integration/integrate-special/4}{\(1/\log x\) is the logarithmic integral,
returned as \(\operatorname{Ei}(\log x)\).}
\pair{integration/integrate-special/5}{\(\sin x/x\) gives the sine integral
\B{SinIntegral}.}
\end{codepairs}

Pairs 2 to 5 are not approximations or placeholders. Each special function is
\emph{defined} as the integral in question, and each appears because the
elementary functions have been shown not to contain an antiderivative. To make
``elementary'' precise we need a little algebra.

\begin{definition}
A \emph{differential field}\index{differential field} is a field \(K\) with a
derivation \(D\colon K\to K\), a map that satisfies \(D(a+b)=Da+Db\) and
\(D(ab)=a\,Db+b\,Da\). Its \emph{constants} are the elements with \(Dc=0\). The
\emph{elementary extensions}\index{elementary extension} of \(K\) are built by
adjoining, one at a time, elements \(\theta\) of three kinds:
\begin{itemize}
\item a \emph{logarithm}: \(D\theta = Du/u\) for some \(u\) in the field so far;
\item an \emph{exponential}: \(D\theta = \theta\,Du\);
\item an \emph{algebraic} element: \(P(\theta)=0\) for a polynomial \(P\) with
coefficients in the field so far.
\end{itemize}
A function is \emph{elementary} if it lies in an elementary extension of
\(\mathbb{C}(x)\) with \(D=d/dx\).
\end{definition}

The trigonometric and inverse trigonometric functions are elementary in this
sense, because \(\sin x = (e^{ix}-e^{-ix})/(2i)\) and \(\arctan x =
\tfrac{i}{2}\log\frac{1-ix}{1+ix}\). So are all the radicals. The adjoined
elements are called \emph{monomials}\index{monomial (differential algebra)}, and
the chain of fields \(\mathbb{Q}(x)=K_0\subset K_1\subset\cdots\subset K_n\) is a
\emph{tower}\index{tower (differential fields)}.

\begin{theorem}[Liouville, 1835; Rosenlicht, 1968]
\label{thm:liouville}
\index{Liouville's theorem}
Let \(K\) be a differential field of characteristic zero, and let \(f\in K\). If
\(f\) has an antiderivative in some elementary extension of \(K\) with the same
constants, then there are constants \(c_1,\dots,c_m\), algebraic over the constants of
\(K\), and elements \(v, u_1,\dots,u_m\) of \(K(c_1,\dots,c_m)\), such that
\[
  f \;=\; Dv + \sum_{i=1}^{m} c_i\,\frac{Du_i}{u_i},
  \qquad\text{that is,}\qquad
  \int f \;=\; v + \sum_{i=1}^{m} c_i \log u_i .
\]
\end{theorem}

The theorem is the reason integration is decidable. An elementary antiderivative,
if one exists, cannot be anything: it is an element of the integrand's own field,
plus a finite sum of constant multiples of logarithms of elements of that field.
New exponentials, new radicals and nested logarithms never help. The algorithms of
this chapter are all ways of searching that restricted space in finitely many
steps, so that an empty search is a proof. The proofs Mathilda uses come from
Bronstein's \emph{Symbolic Integration~I}~\cite{bronstein2005} for the
transcendental case, from Trager's thesis~\cite{trager1984} for the algebraic
case, and from Bronstein's 1990 paper~\cite{bronstein1990} for the mixed case.

The decision is exposed as a predicate of its own:

\begin{codepairs}
\pair{integration/elementary-q/1}{\(e^x/x\) is \emph{not} elementary. This is a
proof, not a timeout.}
\pair{integration/elementary-q/2}{Nor is \(e^{x^2}\).}
\pair{integration/elementary-q/4}{But \(e^x/(1+e^x)\) is: its integral is
\(\log(1+e^x)\).}
\pair{integration/elementary-q/5}{And \(\sin x/x\) is not. The sine integral
exists because it is not.}
\end{codepairs}

\mcode{Risch\textasciigrave ElementaryIntegralQ} has three answers.
\texttt{True} means an elementary antiderivative has been \emph{exhibited} and
verified by differentiation. \texttt{False} means a theorem has been applied: the
decision procedure found an exact certificate that no antiderivative of the
Liouville form exists. Anything else leaves the call unevaluated, with a
\texttt{::undec} message. Section~\ref{ssec:int-verdicts} returns to this.

%======================================================================
\section{Rational functions}
\label{sec:int-rational}

Rational functions are the base of every tower, and for them integration is
completely solved.

\begin{codepairs}
\pair{integration/integrate-elementary/1}{Term-by-term integration of a
polynomial.}
\pair{integration/integrate-elementary/2}{A logarithmic derivative recognised
directly: \(\int 2x/(x^2+1) = \log(1+x^2)\).}
\pair{integration/integrate-elementary/3}{The arctangent, from an irreducible
quadratic denominator.}
\pair{integration/integrate-elementary/4}{A partial-fraction case. The answer
combines two of the three logarithms into a single quadratic \(\log\).}
\pair{integration/integrate-elementary/5}{Two irreducible quadratic factors give a
mixed \(\log\)/\(\arctan\) answer.}
\end{codepairs}

The algorithm has two stages. \emph{Hermite reduction}\index{Hermite reduction}
removes the repeated factors of the denominator. Write \(f=A/D\) and let
\(D=D_1D_2^2\cdots D_k^k\) be its squarefree factorization. Repeatedly solving a
linear congruence with the extended Euclidean algorithm writes
\[
  \int \frac{A}{D} \;=\; \frac{B}{E} + \int \frac{C}{D^{*}},
\]
where \(D^{*}=D_1D_2\cdots D_k\) is squarefree. The rational part \(B/E\) is
found by linear algebra alone, with no factoring. What is left has only simple
poles, so its integral is a sum of logarithms. Its coefficients are the residues
\(c=C(a)/D^{*\prime}(a)\) at the roots \(a\) of \(D^{*}\). Rothstein and Trager
observed that the residues are exactly the roots of the resultant
\[
  R(z) \;=\; \operatorname{Res}_x\bigl(C - z\,D^{*\prime},\; D^{*}\bigr),
\]
and that the logarithm belonging to a residue \(c\) is
\(\log\gcd(C-c\,D^{*\prime}, D^{*})\). Lazard, Rioboo and Trager showed how to
read those gcds from the subresultant sequence of the same resultant computation,
so the splitting field of \(D^{*}\) is never needed. Only the residues need
algebraic numbers, and then only the roots of the factors of \(R\).
\calloutnote{Under the hood}{The rational integrator is
\texttt{src/calculus/intrat.c}, called by the cascade as the
\texttt{"BronsteinRational"} stage. It implements Hermite reduction and
Lazard--Rioboo--Trager\index{Lazard--Rioboo--Trager algorithm}\index{Rothstein--Trager resultant},
never an undetermined-coefficient solve. Rioboo's conversion turns each
complex-conjugate pair of logarithms into a real \(\arctan\), which is why pair~5
comes out in real form. For plain rational input the polynomial arithmetic is
handed to FLINT~\cite{flint}. The classical account is
Geddes--Czapor--Labahn~\cite{gcl} and Chapter~2 of Bronstein's
book~\cite{bronstein2005}. Every branch ends at the same check:
\(D[\text{result}]-f\) must simplify to \(0\).}

Rational functions of a radical come next. The simplest of them are rationalized
by a substitution, and the cascade tries those first:

\begin{codepairs}
\pair{integration/integrate-elementary/6}{A rational function of \(\sqrt x\),
rationalized by \(u=\sqrt x\).}
\pair{integration/integrate-elementary/7}{A square root of a quadratic, handled by
an Euler substitution.}
\end{codepairs}

A substitution works when the curve \(y^2=x^2-1\) can be parametrized by rational
functions, that is, when it has genus~\(0\). Section~\ref{sec:int-trager} deals with
the curves that cannot.

%======================================================================
\section{The transcendental tower}
\label{sec:int-tower}

\subsection{Logarithms, exponentials and tangents}
\label{ssec:int-levels}

A transcendental integrand is first written in a tower. Mathilda reads the
kernels \(\log u\), \(e^u\) and \(\tan u\) out of the integrand, orders them so that
each one's argument lies in the field below it, and integrates from the top down.
At the top level the integrand is a rational function of one monomial \(\theta\)
over the field \(k\) below, and the algorithm is a version of the rational one:
Hermite reduction in \(\theta\), a resultant for the logarithmic part, and then a
\emph{polynomial part} that needs a new tool, depending on the kind of \(\theta\).

\begin{codepairs}
\pair{integration/tower/1}{Three logarithm levels: \(\theta_1=\log x\),
\(\theta_2=\log\theta_1\), and an answer one level higher still.}
\pair{integration/tower/2}{A polynomial in \(\log\log x\). Its coefficients are
integrated in the field below.}
\pair{integration/tower/3}{\(x^x=e^{x\log x}\): an exponential over a logarithm.}
\pair{integration/tower/4}{A rational function of \(e^x\). The residues are
\(\pm i/2\), and the conjugate logarithms combine into a real \(\arctan\).}
\pair{integration/tower/5}{A \emph{hypertangent} level: \(\theta=\tan x\) with
\(D\theta = 1+\theta^2\).}
\pair{integration/tower/6}{A polynomial in \(\tan x\), reduced degree by
degree.}
\end{codepairs}

For a logarithm, \(D\theta\) lies in \(k\), and the polynomial part
\(\sum a_i\theta^i\) is integrated coefficient by coefficient from the top. Each
coefficient's integral must lie in \(k\), except that it may produce one new
multiple of \(\theta\) itself. This is \emph{limited integration}\index{limited
integration}: integrate \(a\) in \(k\) modulo a given finite set of logarithmic
derivatives. For an exponential \(\theta = e^u\), \(D(\theta^i)=i\,Du\,\theta^i\),
so \(\int a_i\theta^i = y_i\theta^i\) requires
\[
  Dy_i + i\,Du\;y_i \;=\; a_i ,
\]
with \(y_i\in k\). That is the \emph{Risch differential equation}, the heart of the
algorithm and the subject of \S\ref{ssec:int-rde}.

Tangents became monomials of Mathilda's tower in v0.6601. Since
\(\tan u\) is elementary, it could always be rewritten with \(e^{iu}\), and for a
long time that was the only route. But the rewritten integrand often leaves any
tower: \(e^{\tan u}\) becomes the exponential of a rational function of \(e^{2iu}\).
So \(\theta=\tan u\) is now a monomial in its own right, a
\emph{hypertangent}\index{hypertangent monomial}, with \(D\theta = Du\,(1+\theta^2)\).
\calloutnote{Theory}{The hypertangent case of the Risch differential equation is
Bronstein's \(\delta=2\) case~\cite[Ch.~6]{bronstein2005}: special
denominators at the factor \(\theta^2+1\), a degree bound for a nonlinear
derivation, and a cancellation case solved through \(k(i)\). Mathilda adds a
change of variable that links it back to the exponential case. With
\(\alpha=i\) for \(\tan\) and \(\alpha=1\) for \(\tanh\), put
\[
  \eta \;=\; \frac{\alpha+\theta}{\alpha-\theta}.
\]
Then \(\eta\) is the exponential \(e^{-2iu}\) or \(e^{2u}\), and
\(\theta=\alpha(\eta-1)/(\eta+1)\). This Moebius map\index{Moebius change of
variable} is an isomorphism of differential fields over \(k(\alpha)\), so a
``no solution'' found in the exponential field is a theorem about the tangent field
as well. Liouville's theorem allows algebraic constants, so integrating a real
\(\tan\) over \(k(i)\) loses nothing. The one exception is a second tangent level:
\(\tan 2u\) is rational in \(\tan u\), so two tangents can be dependent, and that
case is left undecided.}

\subsection{Liouvillian levels}
\label{ssec:int-liouvillian}

The error function and the exponential and logarithmic integrals are not
elementary, but they are \emph{Liouvillian}\index{Liouvillian monomial}: each is
the integral of an element of the field below. \(\operatorname{erf}u\) has
derivative \(\tfrac{2}{\sqrt\pi}\,Du\,e^{-u^2}\), \(\operatorname{Ei}u\) has
\(Du\,e^u/u\), and \(\operatorname{li}u\) has \(Du/\log u\). So each can be adjoined
as a \emph{primitive} monomial over a field that already contains \(e^{-u^2}\),
\(e^u\) or \(\log u\), and integrated like a logarithm.

\begin{codepairs}
\pair{integration/tower/7}{An \B{Erf} level over \(e^{-x^2}\):
\(\int\operatorname{erf}x\,e^{-x^2} = \tfrac{\sqrt\pi}{4}\operatorname{erf}^2 x\).}
\pair{integration/tower/8}{An \(\operatorname{Ei}\) level over \(e^x\), by the
same coefficient matching.}
\pair{integration/tower/9}{The logarithmic integral \B{LogIntegral} integrates to
an answer with \(\operatorname{Ei}(2\log x)\) in it.}
\end{codepairs}

A level has to be proved \emph{independent} before it is trusted.
\(\operatorname{erfc}u=1-\operatorname{erf}u\), and
\(\operatorname{li}u=\operatorname{Ei}(\log u)\). A tower that held both \B{Erf}
and \B{Erfc} of the same argument as separate monomials would miss that they
differ by a constant.\calloutnote{Pitfall}{This happened. Up to v0.6602,
\(e^u(\operatorname{erf}u+\operatorname{erfc}u)\), which is just \(e^u\), was
certified \emph{non-elementary}. The fix (v0.6603) stopped all certificates over
Liouvillian levels until a proof of independence could be added; v0.6621 then
added that proof for each level. \S\ref{ssec:int-structure} shows the case now
decided correctly. It is the reason every Liouvillian level is now proved
independent before a verdict over it is trusted.}

The same section of the cascade handles the polylogarithm answers. They come from
integrands with a logarithm in a rational function, and they are not decided by
the tower; a separate stage recognises them.

\begin{codepairs}
\pair{integration/integrate-special/6}{A dilogarithm\index{dilogarithm} appears:
\(\int\log x/(1+x) = \log x\log(1+x) + \operatorname{Li}_2(-x)\).}
\pair{integration/integrate-special/7}{A polylogarithm ladder\index{polylogarithm
ladder} for \(x/(e^x-1)\), the Bose--Einstein integrand.}
\end{codepairs}

\subsection{The constant field}
\label{ssec:int-constants}

The constants of the tower matter. Over \(\mathbb{Q}\) the residues of the
logarithmic part are algebraic numbers, and the answer must be written in a field
that holds them. Mathilda's transcendental integrator works over \(\mathbb{Q}\),
over \(\mathbb{Q}(i)\), over number fields generated by the radicals that appear in
the input, and over fields \(\mathbb{Q}(a,b,\dots)\) of symbolic parameters.

\begin{codepairs}
\pair{integration/constants/1}{\(\mathbb{Q}(i)\): the exponential \(e^{ix}\) is a
monomial with \(D\theta = i\theta\).}
\pair{integration/constants/2}{A number field: \(\sqrt2\) in an exponent. The
second term is decided separately, and is an \(\operatorname{Ei}\).}
\pair{integration/constants/3}{A rational function over \(\mathbb{Q}(\sqrt2)\).}
\pair{integration/constants/4}{A parameter \(a\) in the exponent.}
\pair{integration/constants/5}{And a proof over \(\mathbb{Q}(a)\): \(\log(a+x)/x\)
is not elementary for a generic \(a\). Its integral is a dilogarithm.}
\end{codepairs}

A parameter changes what a verdict means. An antiderivative computed over
\(\mathbb{Q}(a)\) is valid for every value of \(a\) except those that make a
denominator, a leading coefficient or a discriminant vanish. Those exceptions form
a finite set, or more generally a proper Zariski-closed set. A certificate of
non-elementarity is weaker in a precise way: it holds for all values of \(a\)
outside a \emph{countable} union of such sets. The standard example is
\(x^a e^x\). It has no elementary integral over \(\mathbb{Q}(a)\), yet for every
\(a\in\{0,1,2,\dots\}\) integration by parts gives one.
\index{generic certificate}\calloutnote{Under the hood}{With parameters present,
the linear systems of the ansatz are solved over \(\mathbb{Q}(\text{params})\)
(\texttt{rt\_solve\_ansatz}, v0.6589--v0.6592). The unknowns are named, so the
parameters are never solved for. Before that, the solver could ``solve'' for \(a\)
itself, pick a special value, and report the special case as the generic one;
five false certificates were traced to it and fixed in v0.6592. The constants \(\pi\), \(e\) and
\(\log 2\) are never treated as parameters.}

\subsection{The structure theorem}
\label{ssec:int-structure}

Before a kernel becomes a monomial it must be shown to be new. \(\log 2x\) is not
independent of \(\log x\), since it is \(\log 2+\log x\). \(\arctan e^x\) is a
combination of \(\log(1\pm ie^x)\), and \(e^{\tanh x}\) is the exponential of a
rational function of \(e^{2x}\). A tower that treated such kernels as independent
would be wrong in both directions. It could miss an antiderivative, and its
certificates would be about a field that does not exist.

The \emph{Risch structure theorem}\index{Risch structure theorem} settles the
question with linear algebra. A new \(\log u\) is algebraic over the tower exactly
when \(Du/u\) is a \(\mathbb{Q}\)-linear combination of the logarithmic derivatives
already present (the \(Du_j/u_j\) of the logarithms and the \(Dv_j\) of the
exponentials \(e^{v_j}\)). A new \(e^v\) is dependent exactly when \(v\) is such a
combination, up to a constant.

\begin{codepairs}
\pair{integration/structure/1}{\(\log 2x\) is folded into the tower as
\(\log 2+\log x\), so this is a polynomial in a single logarithm.}
\pair{integration/structure/2}{The derivative of \(\arctan^3 e^x\) \dots}
\pair{integration/structure/3}{\dots\ integrates back. The \(\arctan\) is folded
as complex logarithms of \(e^x\).}
\pair{integration/structure/4}{\(e^{\tanh x}\): the dependent exponential becomes a
member of the tower, and the verdict is a proof.}
\pair{integration/structure/5}{And the case from the Pitfall above:
\(e^x(\operatorname{erf}x+\operatorname{erfc}x)=e^x\) is now decided
correctly.}
\end{codepairs}

\calloutnote{Under the hood}{The structure theorem is applied when the tower is
built (\texttt{risch\_tower.c}). A dependent logarithm becomes the member
\(\sum r_j t_j + c\), with \(c=(\log\kappa+2\pi i m)/d\) (v0.6609). At most one
transcendental constant (\(\pi\) or one \(\log 2\)) may enter, and before v0.6609
every verdict over such a tower was demoted to undecided. A dependent exponential
becomes \(\prod t_j^{k_j}\) times a constant when the exponents are integers
(v0.6633), and for exponents \(1/m\) the dependency is read the other way
(v0.6644). The relations themselves come from
\mcode{Risch\textasciigrave RationalSpan}, which decides \(\mathbb{Q}\)-linear
relations over algebraic constants; a unique non-rational solution, reduced with
\B{RootReduce}, is a proof.}

\subsection{The Risch differential equation}
\label{ssec:int-rde}

At an exponential level the whole question comes down to
\[
  Dy + f\,y \;=\; g, \qquad f, g\in k,
\]
for an unknown \(y\in k\). Mathilda exposes its solver as
\mcode{Risch\textasciigrave RischDECertified[f, g, x]}. It returns the
solution \(y\), or the string \texttt{"NoSolution"}, or \texttt{\$Failed} when it
cannot decide. \texttt{"NoSolution"} is returned only when every step on the path
to it is a theorem of Chapter~6 of Bronstein's book~\cite{bronstein2005}.

\begin{codepairs}
\pair{integration/rde/1}{\(Dy - y = x\) gives \(y=-1-x\), so
\(\int xe^{-x}=-(1+x)e^{-x}\).}
\pair{integration/rde/2}{\(Dy + 2xy = 1\) has no rational solution. This is the
proof that \(\int e^{x^2}\) is not elementary.}
\pair{integration/rde/3}{\(Dy + y = 1/x\) has none either: \(\int e^x/x\) is not
elementary.}
\pair{integration/rde/4}{The \emph{parametric} equation: for which constant \(c\)
is \(e^x/x + c\,e^x/x^2\) of the form \(Dy+y\)? For \(c=-1/2\), with
\(y=e^x/(2x)\).}
\pair{integration/rde/5}{No constant \(c\) makes \(1/\log x + c/x\) a derivative in
\(\mathbb{Q}(x,\log x)\).}
\end{codepairs}

\calloutnote{Theory}{The solver follows Bronstein~\cite[Ch.~6]{bronstein2005}.
\emph{Weak normalization} and the \emph{normal denominator} bound the denominator
of \(y\) from the poles of \(f\) and \(g\). At \emph{special} polynomials, the
factors of \(\theta\) itself for an exponential or of \(\theta^2+1\) for a tangent,
the power is bounded separately. A degree bound then limits the numerator. The
remaining polynomial equation \(a\,Dq+b\,q=c\) is shrunk by Rothstein's \emph{SPDE}
(\emph{special polynomial differential equation}) reduction\index{SPDE reduction},
which lowers the degree bound one division at a time instead of solving one large
linear system, and the last step is a non-cancellation or cancellation solve.
\emph{Limited integration} and the \emph{parametric} Risch equation
\(Dy+fy=\sum_i c_i g_i\) are what the polynomial parts of logarithmic and
Liouvillian levels reduce to: integrate in \(k\) up to an unknown constant
combination of given functions. \mcode{Risch\textasciigrave ParametricRischDE}
(v0.6620) returns \(\{y,\{c_1,\dots\}\}\) or a proof that no constants exist.}

\subsection{Verdicts}
\label{ssec:int-verdicts}

The decision is three-valued, and the three values are not symmetric.
\texttt{True} needs an antiderivative \emph{exhibited} and verified;
\texttt{False} needs a theorem. Everything else is ``undecided.''

\begin{codepairs}
\pair{integration/verdicts/1}{\(\int e^{\tan x}\) is not elementary, decided over
the hypertangent level.}
\pair{integration/verdicts/2}{But \(e^{\tan u}\sec^2u\) integrates to
\(e^{\tan u}\).}
\pair{integration/verdicts/3}{\(e^x\tan x\): not elementary. The Moebius map turns
the tangent into an exponential over \(\mathbb{Q}(i)\).}
\pair{integration/verdicts/4}{A Liouvillian level: \(\operatorname{erf}x/x\) has
no elementary integral.}
\pair{integration/verdicts/5}{\(e^{-x^2}\operatorname{erf}x\) does, as
\S\ref{ssec:int-liouvillian} showed.}
\pair{integration/verdicts/6}{A radical of \(x\) is outside the transcendental
decision, so the call stays unevaluated (with a \texttt{::undec} message) rather than
guess.}
\end{codepairs}

A \texttt{False} is backed by one of four certificates: a non-constant residue in
the simple part (Bronstein's Theorem~5.6.1), a Risch differential equation with no
solution, a primitive level whose constant term cannot be matched, or the
dilogarithm obstruction in a polynomial in \(\log\). \texttt{True} was once
believed correct ``by construction'' and is now checked: since v0.6589 it is given
only for an antiderivative that differentiates back. The predicate is
transcendental-only. Algebraic integrands are decided by \B{Integrate} itself, with
\texttt{Method -> "Trager"} or \texttt{"Bronstein"}, as the next two sections
show.

Soundness has a cost in presentation, which is worth seeing.

\begin{codepairs}
\pair{integration/elementary-q/6}{A correct answer in an awkward form:
\(\int\cos(\log x)\) is real, but printed with complex powers of \(x\).}
\end{codepairs}

\calloutnote{Pitfall}{The output of pair~6 is exactly
\(\tfrac{x}{2}(\cos\log x + \sin\log x)\). Mathilda returns it as
\((\tfrac14+\tfrac{i}4)x^{1-i}+(\tfrac14-\tfrac{i}4)x^{1+i}\), a real quantity
written with conjugate complex exponents. The transcendental integrator reaches
trigonometric integrands by exponentializing, \(\cos\log x =
(x^{i}+x^{-i})/2\), and the simplifier does not fold the result back into real
form. The value is right, and \B{N} confirms it; the form is a known cosmetic
limitation, not an error.}

%======================================================================
\section{Algebraic functions: Trager's algorithm}
\label{sec:int-trager}

An algebraic integrand such as \(1/\sqrt{x^3+1}\) lives on a curve, here
\(y^2=x^3+1\), and its antiderivative, if elementary, lives in the function field
\(\mathbb{Q}(x,y)\) of that curve together with logarithms of elements of it. When
the curve has genus~\(0\) it can be parametrized and the integral rationalized, as
in \S\ref{sec:int-rational}. When the genus is positive it cannot, and the problem
becomes one about divisors on the curve. Trager's 1984 thesis~\cite{trager1984}
gave the first complete algorithm for it. Mathilda implements it in
\texttt{src/internal/trager/}, one file per chapter of the thesis.

\begin{codepairs}
\pair{integration/trager/1}{\(y^4=x^4+1\) is the Fermat quartic, of genus~3. The
integral is elementary all the same.}
\pair{integration/trager/2}{A cube root on a genus-1 curve, with an
\(\arctan\) and two logarithms.}
\pair{integration/trager/3}{\(\sqrt{x^4+1}\) with a rational factor: one real
\(\arctan\).}
\pair{integration/trager/4}{A \emph{pseudo-elliptic} integral: an elliptic curve,
yet an elementary answer, with a degree-8 logand.}
\pair{integration/trager/5}{Two different radicals. Each component is integrated
on its own curve.}
\pair{integration/trager/6}{Chebyshev's binomial \(\sqrt x\,(x+1)^{1/3}\) \dots}
\pair{integration/trager/7}{\dots\ is not elementary, and the message is a
proof.}
\end{codepairs}

\subsection{The stages}
\label{ssec:int-trager-stages}
\index{Trager's algorithm!stages}

The algorithm takes a rational function of \(x\) and one radical
\(y=r(x)^{1/n}\). Powers of the same radicand are merged first:
\(\sqrt{x}\) and \(x^{1/3}\) are both powers of \(x^{1/6}\).
\begin{enumerate}
\item \emph{Canonical form.} Clear the radicand's denominator, \(\hat
y=q\,y\), so that \(\hat y^n\) is a polynomial, and write the integrand as
\(f=\sum_{i=0}^{n-1}a_i(x)\,\hat y^{\,i}\). This is an exact identity in the field;
no branch of the radical is chosen.
\item \emph{Components.} The term \(a_0\) is rational and goes to the rational
integrator. Each other term \(a_i\hat y^{\,i}\) is integrated \emph{on its own
curve} \(Y^m=G\), where \(Y\) is \(\hat y^{\,i}\) with its rational part removed and
\(m=n/\gcd(n,i)\). The integrand is elementary if and only if every component is
(Trager, Appendix~A). This is why pair~5 can split into an \(\arctan\) from the
square root and an \(\arctan\) and logarithms from the cube root.
\item \emph{Integral basis.} For a curve \(Y^m=G\) the integral basis has a closed
form, \(w_i=Y^i/d_i\) with \(d_i=\prod p_j^{\lfloor i e_j/m\rfloor}\) over the
factorization \(G=\prod p_j^{e_j}\). It is diagonal for the derivation and normal
at infinity, so no general integral-basis algorithm is needed.
\item \emph{Hermite reduction} in that basis. The congruences decouple coordinate
by coordinate, and the algebraic part of the answer comes out by linear algebra,
as in the rational case.
\item \emph{Residues.} If a pole at infinity survives Hermite reduction, the
integral is not elementary. Otherwise the residues at the finite poles are the
roots of a resultant, the closed form of Trager's double resultant.
\item \emph{Divisors.} For each residue the places where the differential has that
residue form a divisor \(\delta\) on the curve, generated as an ideal of the
integral closure.
\item \emph{Torsion.} The logarithmic part exists if and only if some multiple
\(N\delta\) is \emph{principal}, the divisor of a function \(u\); then the
contribution is \(\tfrac{1}{N}\log u\). Deciding whether such an \(N\) exists is the
hard step (below).
\item \emph{Log part and check.} Build the logarithms, subtract their derivatives,
and verify that nothing is left. With Trager's construction a nonzero remainder is
a proof of non-elementarity.
\item \emph{Real form.} \texttt{LogToArcTan} rewrites conjugate complex logarithms
as \(\tfrac{c_1}{2}\log(A^2+B^2) - c_2\arctan(B/A)\), and keeps the complex form if
the real one fails its check.
\end{enumerate}
The last guard is \(D[\text{result}]-f=0\), tested with \B{PossibleZeroQ}. A
failure there is an internal error. It produces a decline, never an answer, so a
bug cannot become a wrong answer or a false certificate.

\subsection{Torsion, one Galois orbit at a time}
\label{ssec:int-torsion}
\index{Trager's algorithm!torsion}

The divisor \(\delta\) of degree zero defines a point of the curve's Jacobian, an
abelian group. The question in step~7 is whether that point has finite order. On
an elliptic curve, \(1/\sqrt{x^3+1}\) gives a point of infinite order, and its
integral is a genuine elliptic integral. For pair~4 the point is torsion, and that
is what makes the integral \emph{pseudo-elliptic}\index{pseudo-elliptic integral}:
the logand is the function whose divisor is the corresponding multiple.

Trager decides finiteness by reduction modulo primes. Over a finite field
\(\mathbb{F}_p\) the Jacobian is a finite group, with at most \((1+\sqrt p)^{2g}\)
elements by the Weil bound, so the order \(o_p\) of the reduced point can be found
by trying \(M=1,2,\dots\). Reduction is injective on torsion of order prime to
\(p\). So take two good primes \(p\) and \(q\). If the parts of \(o_p\) and \(o_q\)
prime to \(pq\) differ, the point cannot have finite order, and the integral is
not elementary. If they agree, there is exactly one candidate \(N\), and one
principality test decides.
\calloutnote{Under the hood}{\texttt{TorsionModP.m} accepts a prime only if the
reduction is \emph{good}: \(p\) divides neither \(m\) nor any denominator, the
squarefree part of \(G\) keeps its degree (so the integral basis and the genus are
preserved), and the residues stay distinct and nonzero. The genus comes from the
Riemann--Hurwitz formula (\texttt{Genus.m}), and divisor classes are added in
reduced form, read off a Popov form of the module (\texttt{Jacobian.m}). The
principality test (\texttt{PrincipalGen.m}) is exact: \(\delta\) is principal if
and only if some row of a basis normal at infinity is regular at infinity.}

The residues need not be rational. Their minimal polynomial over \(\mathbb{Q}\)
splits them into Galois orbits, and Mathilda builds the logarithmic part one orbit
at a time, as a \B{RootSum} over each irreducible factor. In genus~\(0\) this always
succeeds. In higher genus an orbit's divisor can have infinite order while some
\(\mathbb{Q}\)-linear combination across orbits is torsion. That failure is not a
certificate, and the code then falls back to Trager's full construction over a
\(\mathbb{Q}\)-basis of the residues.

The Chebyshev example of pairs 6 and 7 is a classical theorem, and a good test.

\begin{theorem}[Chebyshev, 1853]
\index{Chebyshev's theorem (binomial differentials)}
For rational \(m,n,p\) and nonzero \(a,b\), the binomial differential
\(x^m(a+bx^n)^p\,dx\) has an elementary integral if and only if one of \(p\),
\(\tfrac{m+1}{n}\), \(\tfrac{m+1}{n}+p\) is an integer.
\end{theorem}

For \(\sqrt x\,(x+1)^{1/3}\) the three numbers are \(\tfrac13\), \(\tfrac32\) and
\(\tfrac{11}{6}\), so no elementary integral exists. Mathilda does not consult the
theorem here: \texttt{Method -> "Trager"} reaches the same verdict from the
divisor computation.

\subsection{Constants and parameters}
\label{ssec:int-trager-fields}

The constant field can be \(\mathbb{Q}(i)\), a number field generated by radicals of
rationals in the input, or a field of parameters \(\mathbb{Q}(a,b,\dots)\).

\begin{codepairs}
\pair{integration/trager-fields/1}{Over \(\mathbb{Q}(i)\): reduction then uses
primes \(p\equiv1\pmod 4\), where \(-1\) has a square root for \(i\) to map
to.}
\pair{integration/trager-fields/2}{Over \(\mathbb{Q}(\sqrt2)\), with the radical
constant kept literal.}
\pair{integration/trager-fields/3}{A pole at the irrational point \(x=\sqrt2\).}
\pair{integration/trager-fields/4}{A parameter: the answer holds for every
\(a\neq0\).}
\pair{integration/trager-fields/5}{One more \(x\) term in the radicand, and the
integral is not elementary for generic \(a\).}
\pair{integration/trager-fields/6}{\(1/\sqrt{x^3+a}\): an elliptic integral over
\(\mathbb{Q}(a)\), certified.}
\pair{integration/trager-fields/7}{Two parameters, and a derivative recognised.}
\end{codepairs}

\calloutnote{Theory}{Over \(\mathbb{Q}(a)\) the torsion test cannot reduce modulo a
prime directly, so Mathilda first \emph{specializes} the parameter to a good
rational value \(a_0\). A value is good when it preserves the degree, genus and
squarefree profile of the curve, the degrees of the places, and the
irreducibility of each residue factor. Specialization is injective on torsion, so
if the specialized point has infinite order, so has the generic one; otherwise
the generic order is the specialized order \(N_0\) exactly when \(N_0\delta\) is
principal over \(\mathbb{Q}(a)\). Hilbert's irreducibility theorem guarantees that
good values exist. The verdicts carry the generic meaning of
\S\ref{ssec:int-constants}.}

\subsection{Choosing the method}
\label{ssec:int-trager-method}
\index{Trager method}

\texttt{Method -> "Trager"} runs the algorithm alone, with no time budget.
\texttt{Method -> "TragerKauers"} replaces the divisor computation in the
logarithmic part with Kauers' heuristic~\cite{kauers2008}, which uses Gr\"obner bases to look
for the least \(k\) for which the \(k\)-th power of an ideal attached to the
residues is principal. It is often
as fast, but it is a heuristic: when it does not close, it declines, and only the
shared stages can give it a certificate.
\calloutnote{Performance}{The stress report \texttt{tests/data/trager\_stress\_report.md}
(2026-11-22) runs 378 algebraic integrands through both methods and FriCAS, with a
60~s limit. \texttt{"Trager"} agrees with FriCAS on 180 elementary and 132
non-elementary cases, decides 3 that FriCAS does not, and has no case FriCAS
decides that it misses, no disagreement, no wrong answer and no timeout. 63 of its
certificates are ``unconfirmed'' only because FriCAS errored or timed out on them.
\texttt{"TragerKauers"} misses 55 cases FriCAS decides. Trager's median time is
78~ms, against FriCAS's 120~ms.}

%======================================================================
\section{Mixed integrands: Bronstein's algorithm}
\label{sec:int-bronstein}

A radical can sit on top of a tower of logarithms and exponentials:
\(\sqrt{x+\log x}\) is algebraic over \(\mathbb{Q}(x,\log x)\). Bronstein's 1990
paper~\cite{bronstein1990} extends Trager's algorithm to such integrands, and
Mathilda implements it, with extensions, in \texttt{src/internal/bronstein/}.

\begin{codepairs}
\pair{integration/bronstein/1}{An elementary case: a derivative recognised over
\(\mathbb{Q}(x,\log x)(y)\), \(y^2=x+\log x\).}
\pair{integration/bronstein/2}{\(\int\sqrt{x+\log x}\) \dots}
\pair{integration/bronstein/3}{\dots\ is not elementary. \texttt{Last[\$MessageList]}
names the message that \B{Check} caught.}
\pair{integration/bronstein/4}{A radical over an exponential.}
\pair{integration/bronstein/5}{Two radicals and an exponential \dots}
\pair{integration/bronstein/6}{\dots\ and the integral of the derivative is the
function back.}
\pair{integration/bronstein/7}{A radical at the bottom, a logarithm above it: no
elementary integral.}
\pair{integration/bronstein/8}{A kernel whose argument contains the radical: no
elementary integral either.}
\end{codepairs}

There are two shapes. In \emph{shape~A}, ``radical on top,'' the field is
\(k(\theta)(y)\) with \(y^n=P\in k[\theta]\), where \(k\) is a tower and \(\theta\)
its top monomial, as in pairs 1 to 3. In \emph{shape~B}\index{Bronstein's
algorithm!shape B}, ``radical at the bottom,'' the field is
\(\mathbb{Q}(x)(y)(\theta_1,\dots,\theta_r)\) with \(y^n=g(x)\): transcendental
Risch over an algebraic base, as in pairs 7 and 8.

\calloutnote{Under the hood}{Shape~A follows Bronstein's pipeline: Hermite
reduction at the normal places, a separate reduction at the \emph{special}
places (the zeros and the pole of \(\theta\) for an exponential), the resultant
criterion for the residues (a non-constant residue is a certificate), divisors,
torsion, and the logarithmic part. Two choices simplify it. Each component
\(c\,Y\) is handled on its own, as in Trager, and its algebraic part is then the
Risch differential equation \(Db+sb=\tilde c\) over \(k(\theta)\) with
\(s=DG/(mG)\), solved by \mcode{Risch\textasciigrave RischDECertified}. That
replaces Bronstein's Puiseux-based systems for the algebraic part. At the
special places, the binomial curve needs no Newton polygon: the Puiseux
expansions\index{Puiseux expansion} have the closed form
\(Y=\zeta\gamma\,\theta^{g/m}\hat G^{1/m}\), and their coefficients come from
J.~C.~P.~Miller's recurrence. Torsion is decided by specializing \(x\) and the
lower monomials to rationals and then reducing modulo~\(p\) as in Trager. That
step is new: Bronstein notes~\cite[p.~130]{bronstein1990} that the Weil bound does
not apply over a transcendental \(k\).}

\subsection{Risch over the curve}
\label{ssec:int-phase5}
\index{Bronstein's algorithm!over an algebraic base}

Shape~B with a kernel whose argument is free of \(y\) splits by components: every
Risch differential equation it needs is over the transcendental field \(k\), and
\texttt{RDEAlgebraic.m} solves the rest by bounds and undetermined coefficients
over \(\mathbb{Q}(x)(y)\), where ``no solution'' is again a proof. The general
case, a radical anywhere in the tower, needs the Risch algorithm over the curve
itself: a differential field \(\mathbb{Q}(x)(y)(t_1,\dots)\) whose base is an
algebraic function field, with the Risch differential equation solved on the curve.
This is the project's ``Phase~5.''

The theory note \texttt{docs/notes/bronstein\_extensions.tex} (sections E1 and E2)
works out what that requires. Three of its ingredients do not appear in
Bronstein's paper~\cite{bronstein1990}, his paper on the Risch equation on a
curve~\cite{bronstein1991}, or his book~\cite{bronstein2005}: oracles for
logarithmic derivatives over a curve, a degree bound that does not assume
normality at infinity, and an admission rule that keeps the constants in
\(\mathbb{Q}\). The note also observes that the base field \(\mathbb{Q}(x)(y)\) is
not worked out in those sources. One kernel whose argument contains \(y\), as in
pair~8, was already handled (\texttt{BaseCurve.m}). The first milestone of the
general case handles at most two monomials over \(\mathbb{Q}(x)(y)\), and decided
two class-B cases of the integration shootout (\S\ref{ssec:int-fricas}) that no
other stage did. More monomials, and a radical in the
middle of the tower, are still open.

\subsection{Several radicals}
\label{ssec:int-galois}

Two independent square roots \(y_1, y_2\) generate a field with Galois group
\(\{\pm1\}^2\) over the base. The projections
\[
  f_e \;=\; 2^{-m}\sum_{\sigma}\chi_e(\sigma)\,\sigma(f)
\]
onto the characters \(\chi_e\) of that group split \(f\) into \emph{Galois
components}\index{Galois components (integration)}, each a single radical
\(\prod y_i^{e_i}\) times an element of the base. Since \(f\) is elementary if and
only if every component is, each can go to the one-radical algorithm.
When the radicals are not independent in that way, Mathilda groups the terms of the
integrand by the radicals and kernels they share, into \emph{support
groups}\index{support groups (integration)}. If every group is elementary, the
answers add. If exactly one group is certified non-elementary and the others are
elementary, the whole integrand is non-elementary.

\begin{codepairs}
\pair{integration/bronstein-split/1}{Two square roots, each times \(\log x\):
two components, both elementary.}
\pair{integration/bronstein-split/2}{One support group elementary, the other
\(e^x/x\): so the sum is not elementary.}
\end{codepairs}

With two or more non-elementary groups no conclusion follows. Their
non-elementary parts might cancel, and the separation lemma that would rule this
out is not proved. Mathilda then reports the case as incomplete.
\index{Bronstein method}\texttt{Method -> "Bronstein"} runs this stage alone. It
hands a radical-free integrand to the transcendental integrator and an integrand
with no logarithm or exponential to Trager.

%======================================================================
\section{ParallelMixedTower}
\label{sec:int-pmt}
\index{ParallelMixedTower method}

The decision procedures cover precisely defined classes. Many integrands fall just
outside them: several radicals at once, a nested radical, a radical of a
trigonometric function. The \texttt{"ParallelMixedTower"} method reaches them by a
different route. It is a \emph{parallel} Risch--Norman integrator: it writes the
whole integrand in a mixed tower built from rational operations, radicals, and
\(\log\)/\(\exp\)/\(\tan\) kernels, guesses the shape of the antiderivative from
the denominator and the kernels, and finds the unknown coefficients by one linear
solve. It verifies every result by differentiating it before returning it. A case
it cannot realize correctly declines. It does not return a wrong answer.

\begin{codepairs}
\pair{integration/integrate-special/8}{\(\int\sqrt{\tan x}\,dx\): a radical of a
tangent, closed in real \(\arctan\)/\(\log\) form.}
\pair{integration/pmt/1}{Two radicals, \(\sqrt x\) and \(\sqrt{x+1}\): the
\emph{Kummer split}.}
\pair{integration/pmt/2}{A conic with two parameters, parametrized from a point
over \(\mathbb{Q}(a,b)\).}
\pair{integration/pmt/3}{A radical of \(\log x\): the substitution \(x=e^u\).}
\pair{integration/pmt/4}{A radical of \(e^x\): the substitution \(u=e^x\).}
\pair{integration/pmt/5}{A nested radical.}
\end{codepairs}

Three preprocessing steps, added in v0.6099--v0.6112, account for most of what the
method gains on hard corpora.
\begin{itemize}
\item The \emph{Kummer split}\index{Kummer split} (v0.6099, v0.6103) writes an
integrand that no single tower holds as \(\sum g_j M_j\), where each \(M_j\) is a
product of radicals, groups the terms by Kummer class, and integrates each group
with one radical, trying Trager first. A non-elementary group is a certificate for
the whole integrand.
\item \emph{Kernel substitutions} (v0.6105) are the differential-field
isomorphisms \(u=\sin x\), \(u=\cos x\), \(u=e^x\), \(x=u^k\), \(x=e^u\) (for
radicals of \(\log\)) and \(x=\tan u\) (for radicals of \(\arctan\)), plus an Euler
parametrization for \(\sqrt{Q(x^k)}\).
\item \emph{Conics with parameters} (v0.6106): a conic radicand over
\(\mathbb{Q}(\text{params})\) is parametrized from a point found by
\texttt{ParametricPoint}, and radicals of parameters become new indeterminates. No
certificate is issued over \(\mathbb{Q}(\text{params})\) here; a certificate needs
the generic argument of \S\ref{ssec:int-trager-fields}, which this method does not
make.
\end{itemize}
\calloutnote{Performance}{On \emph{Diabolically Difficult Integrals},
Volume~1, a book of 803 integrands, \texttt{"Trager"} and
\texttt{"ParallelMixedTower"} together decide 664 cases: an antiderivative, or a
proof that none is elementary (v0.6112). FriCAS decides 528 of the 803. The
parallel method alone went from 490 to 564 in the last of those releases, with no
case lost and no wrong answer. It is also slower than the dedicated methods, which
is why it is placed after the decision procedures.}

%======================================================================
\section{The cascade and its guarantees}
\label{sec:int-cascade}
\index{Integrate!cascade}

With no \texttt{Method} option, \B{Integrate} runs a \emph{cascade}: a fixed
sequence of stages, each correct on the family it accepts. The first stage that
returns a verified answer wins. In order, the indefinite cascade runs:
\begin{enumerate}
\item cheap structural stages, the rational integrator, and the classical radical
substitutions (linear, quadratic and linear-ratio radicals, Chebyshev's binomials,
Fresnel and power-of-Gamma forms);
\item \texttt{"Trager"} and then \texttt{"Bronstein"}, each with a budget;
\item the early transcendental Risch stage, for large integrands with an
exponential, logarithmic or trigonometric kernel;
\item linearity, linear-argument substitution, the Weierstrass substitution,
derivative-divides, and the general transcendental Risch stage;
\item the special-function stages: the integral tables, the polylogarithm ladder,
holonomic modules, Bessel and hypergeometric forms, and integration by parts with
logarithms and special functions;
\item \texttt{"ParallelMixedTower"}, then Goursat's pseudo-elliptic forms, then a
special-function search over the mixed tower;
\item Meijer-G and the Gradshteyn--Ryzhik table, last.
\end{enumerate}
\calloutnote{Under the hood}{The order is in \texttt{src/calculus/integrate.c}.
Trager and Bronstein moved ahead of linearity and the transcendental stage in
v0.6521: each of those had been measured grinding for 60~s on integrands Bronstein
closes in under 5~s. A certificate from either stage skips the stages that can
only answer elementarily (ParallelMixedTower and Goursat), since they can no longer
succeed, but not the special-function stages, since an \(\operatorname{erf}\) or
elliptic form may still exist. The variable \texttt{Integrate\textasciigrave
\$CascadeParameters} (default \texttt{\{"Trager"\}}) names the stages that accept
an integrand with symbolic parameters; Bronstein over parameters is opt-in.}

Many of the cascade's stages are small. Four illustrate the range:

\begin{codepairs}
\pair{integration/integrate-newer/1}{A shared linear argument \(3x\) is
substituted away before the Weierstrass stage sees it, so \(\int\sec^2 3x =
\tfrac13\tan 3x\) rather than a high-degree rational function of \(\tan(x/2)\).}
\pair{integration/integrate-newer/2}{A logarithm times a trigonometric kernel, by
one integration by parts. The leftover \(\int\cos x/x\) is the cosine integral
\B{CosIntegral}.}
\pair{integration/integrate-newer/3}{A \emph{holonomic} answer:
\(1/\operatorname{Ai}^2\) integrates to \(\pi\operatorname{Bi}/\operatorname{Ai}\),
by the Wronskian of the two Airy functions.}
\pair{integration/integrate-newer/4}{A Bessel product closed in its own differential
module: \(\int xJ_0^2 = \tfrac{x^2}2(J_0^2+J_1^2)\).}
\end{codepairs}

Pair~2 is the log-by-parts stage\index{integration by parts!logarithmic}. For
\(c\log g(x)\,K(x)\) with a trigonometric or hyperbolic kernel \(K\), it does the
single step \(\int\log g\,K=\log g\,V-\int V\,g'/g\) with \(V=\int K\), and accepts
the result only if it differentiates back exactly. Pairs 3 and 4 are the
\texttt{"Holonomic"}\index{holonomic function} stage. Bessel, Airy, and sine/cosine
kernels of a linear argument each span a two-dimensional space closed under
\(d/dx\), so a power of \(x\) times a product of them lies in a finite-dimensional
\emph{differential module}. One linear system then finds the antiderivative inside
that module, or the stage declines. For \(1/\operatorname{Ai}^2\) another route
applies: with \(t=\operatorname{Bi}/\operatorname{Ai}\), \(t'=W/\operatorname{Ai}^2\)
where the Wronskian \(W=1/\pi\) is constant, so the integral is \(\pi t\).
\calloutnote{Performance}{Before the log-by-parts stage existed, \(\int\log
x\sin x\) fell through to the general \texttt{"ParallelMixedSpecial"} search and
took about \(5\)~s; the stage closes it in roughly a tenth of a second. The saving
matters most inside \B{DSolve}, whose variation-of-parameters quadratures produce
exactly this shape: \(y''+4y=\log x\) dropped from over \(9\)~s to about a third
of a second.}

\subsection{Budgets and messages}
\label{ssec:int-messages}

The two decision stages run under wall-clock budgets, in seconds. When a budget
expires the stage declines, and the cascade continues. \texttt{Infinity} removes a
budget, and an explicit \texttt{Method} runs with none.

\begin{codepairs}
\pair{integration/cascade/1}{\mcode{\$TragerCascadeBudget} and
\mcode{\$BronsteinCascadeBudget}: ten seconds each.}
\pair{integration/cascade/2}{The full cascade proves \(\log x/\sqrt{1+x^3}\) not
elementary.}
\end{codepairs}

The two messages make different promises.
\begin{itemize}
\item \texttt{Integrate::nonelem}\index{Integrate::nonelem@\texttt{Integrate::nonelem}}
(``The integrand \dots\ has no antiderivative elementary in \dots'') is printed only
with a proof: a certificate from the transcendental Risch stage, from Trager, from
Bronstein, or a rigorous one from ParallelMixedTower. It is held until the cascade
ends, and printed only if no later stage finds a special-function answer. It names
the user's integrand once, never an inner integral.
\item \texttt{Integrate::incomplete}\index{Integrate::incomplete@\texttt{Integrate::incomplete}}
(``The integrand \dots\ was not decided in \dots: implementation incomplete
(\dots)'') means the integrand is algebraic over an elementary tower, but outside
what is implemented. Either it is outside one simple radical extension (nested
radicals, several radicals whose Kummer group is not cyclic, a general algebraic
function), or it is inside it on a branch not yet written. The reason in
parentheses says which.
\end{itemize}
A third outcome, a \emph{decline}, prints nothing: an approximate number, a
non-elementary function in the input, or a budget that ran out. The integral is
returned unevaluated.

\calloutnote{Performance}{A proof does not end the cascade. The stages after it
may still find a special-function antiderivative, as one does for
\(1/\sqrt{1-x^4}\) (an \B{EllipticF}), so they cannot simply be skipped, and a
non-elementary integrand pays for whatever they try. The certificate itself is
cheap; the cost is in those later stages, and it depends on the integrand. On this
build \(\int dx/(x\sqrt{x^3+ax+1})\) takes about 16~s in the full cascade against
under 0.1~s with \texttt{Method -> "Trager"}, with the same verdict. Nearly all of
it is ParallelMixedSpecial: it finds a special-function antiderivative quickly, then
spends its time building the elliptic integral of the third kind over the cubic's
irrational roots, until its budget cuts it off. Chebyshev's
\(\sqrt x\,(x+1)^{1/3}\) takes about 7~s against under 0.1~s, spent in
DerivativeDivides (about 1.7~s) and the Roach MeijerG stage (about 5~s).
(\texttt{MATHILDA\_INTEGRATE\_PROFILE=1} in the environment prints this per-stage
breakdown.) That is why the slower examples above that end in a proof name the method. When
you expect a ``no,'' ask the decision procedure directly.}

\subsection{How it compares}
\label{ssec:int-fricas}
\index{FriCAS}

FriCAS, the descendant of Axiom, has the most complete implementation of the
Risch--Trager--Bronstein algorithm in any other system; Bronstein wrote much of it.
Mathilda's methods were measured against FriCAS~1.3.6 on four corpora. Every
answer was verified by the system that gave it.

\begin{center}
\small
\begin{tabular}{@{}llrr@{}}
\hline
Corpus & Cases & Mathilda decides & FriCAS decides\\
\hline
256, integration shootout & 1266 & 96\% (1211) & 75\% (952)\\
257, compact integrands, large answers & 1000 & 100\% (995) & 90\% (899)\\
259, transcendental decision & 640 & 98\% (625) & 95\% (610)\\
260, radicals at the bottom & 420 & 90\% (377) & 95\% (397)\\
\hline
\end{tabular}
\end{center}

``Decides'' means a verified closed form, or a proof that no elementary integral
exists. The first three rows are from the 2026-11-21 report
(\texttt{benchmarks/BRONSTEIN\_VS\_FRICAS.md}, Mathilda v0.6646); the last is from
v0.6648. On algebraic integrands (class~A of corpus 256) Mathilda decides 100\%
against 82\%, and on a radical over a logarithm or exponential (class~M) 100\%
against 44\%. Radicals at the bottom (class~B, and corpus 260) are where FriCAS led,
and where most of the recent work went: class~B rose from 72\% to 168 of 169 by
v0.6648, and by v0.6651 every family of corpus 260 decides at least 56 of its 60
cases. No Mathilda method gave a false ``no'' or a wrong answer on any corpus.
FriCAS gave 14 false ``no''s on corpus 259, all exponentials over
\(\mathbb{Q}(i)\), now reported upstream. FriCAS still \emph{closes} more
transcendental cases with special functions where Mathilda proves
non-elementarity. Mathilda is faster on the typical case (medians of 16--71~ms
against 70--110~ms) but has a slow tail: 51 common cases take over a second, against
none for FriCAS. Most of that tail is the Bronstein package, which is written in
Mathilda's own language.

\subsection{What is still open}
\label{ssec:int-open}

The decision procedures are complete only on the classes they implement. The
transcript below shows where those classes end today.

\begin{codepairs}
\pair{integration/cascade/3}{A radical in the middle of the tower,
\(\log(1+\sqrt{\log x})/x\): the decision procedure reports it incomplete.}
\pair{integration/cascade/4}{The cascade closes this one anyway, by the
substitution \(u=\log x\). Not every case has such a route.}
\pair{integration/cascade/5}{A nested radical: incomplete, for Bronstein's
method.}
\pair{integration/cascade/6}{The full cascade decides it all the same:
ParallelMixedTower's rigorous certificate proves it non-elementary.}
\pair{integration/cascade/7}{\(\int\sqrt{\tan x}\) over \(\mathbb{Q}(i)\) by
\texttt{"Bronstein"} carries \B{Root} objects. (ParallelMixedTower gives the real
form shown in \S\ref{sec:int-pmt}.)}
\end{codepairs}

\begin{tcolorbox}[enhanced, breakable, colback=black!2, colframe=black!45,
  boxrule=0.4pt, left=7pt, right=7pt, top=4pt, bottom=5pt,
  colbacktitle=black!7, coltitle=black, fonttitle=\bfseries,
  title={What is still open (v0.6664)}]
\begin{itemize}
\item \emph{A radical in the middle of a tower.} \(\log(1+\sqrt{\log x})/x\) has a
radical inside a logarithm, with the radicand itself transcendental. Bronstein's
method reports \texttt{Integrate::incomplete}. The general case is the unfinished
part of Phase~5.
\item \emph{Nested radicals.} \(\sqrt{\log x+\sqrt{\log x}}\) is outside one simple
radical extension, and Bronstein's method reports it incomplete; this one the
cascade still proves non-elementary through ParallelMixedTower (pair~6), but the
class is open. Several square roots with two or more non-elementary support groups
are also undecided.
\item \emph{Answers over \(\mathbb{Q}(i)\)} from Trager and Bronstein carry
\B{Root} objects of the residue field, with no radicals and no real form.
\item \emph{\mcode{Risch\textasciigrave ElementaryIntegralQ} is
transcendental-only.} An algebraic integrand such as \(1/\sqrt{1+x^3}\) stays
unevaluated; use \B{Integrate} with \texttt{Method -> "Trager"} or
\texttt{"Bronstein"}. A transcendental constant such as \(\log 2\) beside the tower
demotes a correct ``no'' to undecided.
\item \emph{Real forms.} \(\int\cos(\log x)\) is returned with complex powers of
\(x\) (\S\ref{ssec:int-verdicts}).
\item \emph{Speed.} A non-elementary integrand pays for the special-function stages
after its proof, and the Bronstein package has a slow tail
(\S\ref{ssec:int-fricas}).
\end{itemize}
\end{tcolorbox}

\usagebox{Integrate}

Integration has a second half. A definite integral is not an antiderivative
evaluated at two points when the integrand has poles, the interval is infinite, or
the antiderivative crosses a branch cut; \S\ref{ssec:calc-definite} describes the
definite integrator, which uses residues, Mellin transforms and differentiation
under the integral sign as well as the antiderivatives of this chapter. And the
algorithms here are the quadratures behind \B{DSolve}
(\S\ref{ssec:calc-dsolve}): every integrating factor and every
variation-of-parameters step is a call to \B{Integrate}.
